\documentclass[11pt,a4paper]{article}

\usepackage{fontspec}
\usepackage{babel}
\babelprovide[import, main]{russian}
\babelfont{rm}{Noto Serif}
\babelfont{sf}{Noto Sans}
\babelfont{tt}[AutoFakeBold=1.5,AutoFakeSlant=0.15]{Noto Mono}
\emergencystretch=3em
\tolerance=1400
\hfuzz=0.5pt
\usepackage[a4paper,top=19mm,bottom=19mm,left=21mm,right=21mm,headheight=14pt]{geometry}
\usepackage{amsmath,amssymb,amsthm,mathtools}
\usepackage{booktabs,tabularx}
\usepackage{xcolor}
\usepackage{fancyhdr}
\usepackage{microtype}
\usepackage[colorlinks=true,linkcolor=blue!55!black,citecolor=blue!55!black,urlcolor=blue!55!black]{hyperref}

\definecolor{accent}{RGB}{24,62,112}
\definecolor{muted}{RGB}{95,105,118}
\definecolor{boxbg}{RGB}{245,247,250}

\newtheoremstyle{cleanthm}{8pt}{8pt}{\itshape}{}{\bfseries\color{accent}}{.}{.5em}{}
\newtheoremstyle{cleandef}{8pt}{8pt}{\normalfont}{}{\bfseries\color{accent}}{.}{.5em}{}
\theoremstyle{cleanthm}
\newtheorem{theorem}{Теорема}[section]
\newtheorem{lemma}[theorem]{Лемма}
\newtheorem{proposition}[theorem]{Предложение}
\theoremstyle{cleandef}
\newtheorem{definition}{Определение}[section]
\newtheorem{example}[theorem]{Пример}
\newtheorem{remark}[theorem]{Замечание}

\pagestyle{fancy}
\fancyhf{}
\fancyhead[L]{\small\color{muted}Наблюдать, менять представление или действовать (v2.6.1, испр.)}
\fancyhead[R]{\small\color{muted}А. Странд}
\fancyfoot[C]{\small\thepage}
\renewcommand{\headrulewidth}{0.4pt}
\renewcommand{\footrulewidth}{0pt}

\DeclareMathOperator{\RepairReach}{RepairReach}
\DeclareMathOperator{\obs}{obs}
\DeclareMathOperator{\Gain}{Gain}

\begin{document}

\begin{center}
  {\LARGE\bfseries Наблюдать, менять представление или действовать:\\[3pt]
  сертифицированный выбор стратегии при ограничениях памяти, бюджета и времени}\\[8pt]
  {\large\color{muted} Разделение информационных, структурных и каузально-временных препятствий для ограниченного наблюдателя}\\[10pt]
  {\normalsize Александр Странд}\\[3pt]
  {\small\color{muted} 24 ноября 2026\,г.\,$\cdot$\,препринт v2.6.1, исправленная редакция\,$\cdot$\,формализация, верификация в Lean~4.36.0 (27/27) и Python-эталон}
\end{center}

\vspace{4pt}
\begin{abstract}
\noindent
Представим агента или систему управления, которая должна распознать устройство среды или предотвратить в ней сбой, но заходит в тупик. Причина неудачи может быть одной из четырёх принципиально разных: (1)~не хватает данных --- нужно задать ещё один вопрос в рамках текущей модели; (2)~слишком беден сам язык описания или память --- никакие новые наблюдения не помогут, пока система не перейдёт, например, от стековой памяти к очереди; (3)~слепо окно наблюдения --- два разных состояния мира дают одинаковые показания датчиков, и нужно либо открыть новый канал наблюдения (не перепутав его с бесполезным шумом), либо активно вмешаться в среду; (4)~истекает физическое время --- пока система ждёт новых данных, среда эволюционирует, и окно для безопасного вмешательства закрывается. Обычный алгоритм поиска не различает эти четыре ситуации: при бедности языка он тратит бюджет на бесконечные пробы, при ресурсном таймауте ложно усложняет модель, а при переносе опыта в новую задачу применяет старые инварианты там, где изменился физический фон среды.

В работе построена и машинно верифицирована единая теория того, как ограниченный наблюдатель может строго различать эти препятствия и выбирать допустимое действие с учётом причины тупика. Во-первых, мы формализуем три координаты доступности среды --- выразимость в текущем языке ($\mathcal{X}$), различимость через пассивное окно наблюдения ($\mathcal{D}$) и операционную идентифицируемость через вмешательство ($\mathcal{J}$) --- и доказываем их полное взаимное отделение в единой модели. Во-вторых, мы доказываем безусловное отделение потоковой памяти типа очереди (FIFO) от заданного семейства стековых вычислителей (LIFO) произвольной ёмкости на многосимвольных словах и строим вычислимый детектор мнимых расширений сенсора, отсекающий датчики, не несущие новой разделяющей информации. В-третьих, шесть элементарных шагов --- проба, расширение языка, расширение сенсора, хирургическое вмешательство, стохастическая отбраковка по $e$-значениям и перенос памяти --- объединены в единый контроллер с пятимерным векторным бюджетом и разделением привилегий: таймаут по ресурсам без структурного сертификата неопровержимо блокируется и не может вызвать ложную эскалацию модели. В-четвёртых, на примерах Rule~110 мы связываем перенос памяти и выбор действия с динамикой среды. Проверка переноса принимает токен для модели порта с фиксированными соседями, когда предпосылки переноса заданы в описании среды. Эта модель отлична от наблюдения свободного эфирного кольца. В отдельном примере закрытия каузального горизонта ($t^*=3,\ \Delta t=2$) контроллер выбирает немедленный ремонт вместо ожидания пробы. Все 27~теорем проверены ядром Lean~4.36.0 без внешних библиотек (\texttt{exit 0}, 0~\texttt{sorryAx}, 0~\texttt{native\_decide}: 10~без аксиом и 17~только на стандартных аксиомах \texttt{propext}, \texttt{Classical.choice}, \texttt{Quot.sound}); эталонный пакет на Python~3 воспроизводит конечные проверочные примеры.
\end{abstract}

\section{Постановка задачи: четыре природы тупика при ограниченном наблюдении}

Когда исследовательская программа или автономный контроллер работает в частично наблюдаемой среде при ограниченных ресурсах, любая остановка выглядит снаружи одинаково: текущая цель не достигнута. Однако за этим внешним отказом скрываются четыре ортогональных препятствия, требующих совершенно разных ответных действий~\cite{reiter1987,pearl2009,wolpert2008}:
\begin{enumerate}
  \item \textbf{Информационная неопределённость внутри адекватной модели.} Текущий язык гипотез $H$ уже содержит верное описание, а текущий сенсорный порт $\Pi$ способен различить конкурирующие состояния мира $W$, но нужно выполнить ещё одну пробу $q$, чтобы сузить $W$.
  \item \textbf{Структурная бедность представления ($\neg\mathcal{X}$).} Целевой закон среды вообще не выразим в текущем классе памяти $H$ (например, среда копирует поток по правилу очереди FIFO, тогда как наблюдатель ищет гипотезу среди конечных автоматов или магазинных стеков LIFO). В этом режиме любые новые пробы и любое расширение сенсоров лишь сжигают бюджет: необходим переход на более богатый уровень языка $H\mapsto H'$.
  \item \textbf{Слепота пассивного порта наблюдения ($\neg\mathcal{D}$) или коллизия механизмов ($\neg\mathcal{J}$).} Язык $H$ достаточно богат, но разные миры $w_1,w_2\in W$ с противоположным значением целевого свойства дают тождественные показания на текущем порту $\Pi(w_1)=\Pi(w_2)$. Здесь бесполезно менять язык $H$: требуется либо добавить новый нетривиальный датчик $\Pi\mapsto\Pi'$, либо применить активное хирургическое вмешательство $do(a)$, разделяющее скрытые механизмы.
  \item \textbf{Закрытие каузально-временного горизонта.} В эволюционирующей среде $T:\Omega\to\Omega$ с конечной скоростью распространения взаимодействий ожидание пробы в течение $\Delta t>0$ тактов позволяет полностью узнать мир ($|W'_{q,y}|=1$), но к моменту $t+\Delta t$ зона действия актуатора уже выпадает из обратного светового конуса цели к дедлайну $t^*$. В этом случае стремление <<сначала точно узнать, а потом действовать>> приводит к необратимому сбою: контроллер обязан действовать немедленно на $t=0$.
\end{enumerate}

Цель настоящей работы --- свести диагностику всех четырёх препятствий и правила выбора ответного шага в единую формальную систему с машинно проверяемыми сертификатами в Lean~4.36.0.

\section{Три координаты доступности среды и их полное разделение}

Состояние наблюдателя: $E = \langle W, H, \mathrm{eval}_H, \Pi, A, \mathrm{apply}_A, \Phi^*\rangle$, где $W$ --- конечный список совместимых миров, $H$ --- текущий класс доступных гипотез с семантикой $\mathrm{eval}_H:H\to\Omega\to\mathrm{Bool}$, $\Pi:\Omega\to\mathcal{O}$ --- пассивный порт наблюдения, $A$ --- доступные вмешательства с действием $\mathrm{apply}_A:A\to\Omega\to\Omega$, а $\Phi^*:\Omega\to\mathrm{Bool}$ --- целевой инвариант.

\begin{definition}[Трёхкоординатный профиль эпистемической доступности $\mathcal{E}(E)$]
Определим три булевы координаты состояния $E$:
\begin{align}
  \mathcal{X}_E (\text{Выразимость}) &:\iff \exists h\in H,\ \forall w\in W:\ \mathrm{eval}_H(h,w) = \Phi^*(w), \\
  \mathcal{D}_E (\text{Различимость}) &:\iff \forall w_1,w_2\in W:\ \Phi^*(w_1)\ne\Phi^*(w_2) \implies \Pi(w_1)\ne\Pi(w_2), \\
  \mathcal{J}_E (\text{Идентифицируемость}) &:\iff \mathcal{X}_E \land \mathcal{D}_E \land \Bigl(\forall w_1,w_2\in W:\ \Phi^*(w_1)\ne\Phi^*(w_2) \nonumber\\
  &\qquad\quad \implies \exists a\in A:\ \Pi(\mathrm{apply}_A(a,w_1)) \ne \Pi(\mathrm{apply}_A(a,w_2))\Bigr).
\end{align}
\end{definition}

В предшествующих версиях формализации проверки независимости координат $\mathcal{X}_E$, $\mathcal{D}_E$ и $\mathcal{J}_E$ были разнесены по разным модулям с несовпадающими сигнатурами типов. В едином модуле \texttt{certified\_strategy\_selection\_v2\_6\_1.lean} мы строим четыре состояния $E_1,E_2,E_3,E_4$ над одним и тем же базовым пространством миров $\Omega=\mathrm{Bool}$ и доказываем полную теорему разделения.

\begin{theorem}[Полное разделение и включения трёх координат доступности]
\label{thm:h1_separation}
В единой модели над $\Omega=\mathrm{Bool}$ (\texttt{H1\_three\_\allowbreak coordinate\_\allowbreak strict\_\allowbreak separation}):
\begin{enumerate}
  \item Выполняются строгие включения $\mathcal{J}_E \implies \mathcal{X}_E$ и $\mathcal{J}_E \implies \mathcal{D}_E$ (\texttt{H1\_identifiable\_\allowbreak implies\_\allowbreak expressible}, \texttt{H1\_identifiable\_\allowbreak implies\_\allowbreak distinguishable}).
  \item Выразимость не влечёт различимость ($\mathcal{X}_{E_1} \land \neg\mathcal{D}_{E_1}$, \texttt{H1\_E1\_expressible\_\allowbreak not\_\allowbreak distinguishable}): при богатом языке $H$ и <<слепом>> порте $\Pi(w)=\mathrm{const}$ цель выразима, но миры неразличимы.
  \item Различимость не влечёт выразимость ($\mathcal{D}_{E_2} \land \neg\mathcal{X}_{E_2}$, \texttt{H1\_E2\_distinguishable\_\allowbreak not\_\allowbreak expressible}): точный порт $\Pi(w)=w$ разделяет все миры, но вырожденный класс гипотез $H=\{h_0\}$ с $\mathrm{eval}_H(h_0,w)=\mathsf{false}$ не способен выразить $\Phi^*(w)=w$.
  \item Совместная выразимость и пассивная различимость не влекут операционной идентифицируемости ($\mathcal{X}_{E_3} \land \mathcal{D}_{E_3} \land \neg\mathcal{J}_{E_3}$, \texttt{H1\_E3\_expr\_\allowbreak and\_\allowbreak dist\_\allowbreak not\_\allowbreak identifiable}), если доступные вмешательства разрушают наблюдаемое различие ($\mathrm{apply}_A(a,w)=\mathsf{false}$).
\end{enumerate}
\end{theorem}

\section{Отделение памяти типа очереди от стека произвольной глубины и детектор мнимых сенсоров}

Чтобы переключение языка $H\mapsto H'$ и расширение порта $\Pi\mapsto\Pi'$ опирались на вычислимые сертификаты, а не на эвристику, необходимо закрыть два конструктивных вопроса: (а)~как доказать, что задача действительно требует перехода от стековой памяти (LIFO) к очереди (FIFO) на многосимвольных потоках; (б)~как алгоритмически отличить полезное расширение сенсора от мнимого (шумового или дублирующего).

\subsection{Многосимвольный автомат с очередью против стека произвольной ёмкости}

Рассмотрим трёхбуквенный алфавит $\Sigma=\{\mathtt{0},\mathtt{1},\mathtt{\#}\}$ и потоковые слова вида $u\mathtt{\#}v$, где $u,v\in\{\mathtt{0},\mathtt{1}\}^*$. Определим в чистом ядре Lean~4 два операционных автомата с буфером $\mathrm{buf}\in\mathrm{List\ Bool}$, параметром инверсии бита $\mathrm{inv}\in\mathrm{Bool}$ и ёмкостью буфера $\mathrm{cap}\in\mathbb{N}$:
\begin{itemize}
  \item \textbf{Автомат с очередью $\mathrm{evalQueue}(\mathrm{inv},\mathrm{cap},w)$ (уровень $L_4$, FIFO):} до разделителя $\mathtt{\#}$ записывает биты в \emph{конец} списка $\mathrm{buf}\mapsto\mathrm{buf}\mathbin{+\!+}[b]$ (при $|\mathrm{buf}|<\mathrm{cap}$), а после $\mathtt{\#}$ извлекает первый элемент из \emph{начала} списка ($h::\mathrm{tl}$) и проверяет равенство $b = h \oplus \mathrm{inv}$.
  \item \textbf{Автомат со стеком $\mathrm{evalStack}(\mathrm{inv},\mathrm{cap},w)$ (уровень $L_3$, LIFO):} до разделителя $\mathtt{\#}$ кладёт биты на \emph{вершину} стека $\mathrm{buf}\mapsto b::\mathrm{buf}$, а после $\mathtt{\#}$ снимает бит с той же вершины ($h::\mathrm{tl}$) и проверяет $b = h \oplus \mathrm{inv}$.
\end{itemize}

На паре тестовых слов $w_{00}=\mathtt{00\#00}$ (симметричный повтор) и $w_{01}=\mathtt{01\#01}$ (асимметричное копирование потока FIFO) возникает фундаментальное препятствие для любого стека: после прочтения префикса $\mathtt{01\#}$ на вершине LIFO-стека оказывается бит~$\mathtt{1}$, тогда как первым символом после разделителя в $w_{01}$ идёт~$\mathtt{0}$. Если стек не инвертирует биты ($\mathrm{inv}=\mathsf{false}$), он отвергает $w_{01}$ при любой ёмкости $\mathrm{cap}\in\mathbb{N}$; если же стек инвертирует биты ($\mathrm{inv}=\mathsf{true}$), он отвергает симметричное слово $w_{00}=\mathtt{00\#00}$ при любой ёмкости $\mathrm{cap}\in\mathbb{N}$.

\begin{theorem}[Безусловное отделение очереди от стека произвольной глубины]
\label{thm:h2_queue_vs_stack}
Автомат очереди $\mathrm{evalQueue}(\mathsf{false},3,\cdot)$ принимает потоковые копии $w_{00}=\mathtt{00\#00}$, $w_{01}=\mathtt{01\#01}$, $w_{101}=\mathtt{101\#101}$ и отвергает реверс $w_{\mathrm{rev}}=\mathtt{01\#10}$ (\texttt{H2\_queue\_\allowbreak recognizes\_\allowbreak fifo\_\allowbreak copy}), тогда как для \textbf{любого} булева преобразования $\mathrm{inv}\in\{\mathsf{false},\mathsf{true}\}$ и \textbf{любой} ёмкости стека $\mathrm{cap}\in\mathbb{N}$ (\texttt{H2\_gate13\_\allowbreak queue\_\allowbreak separates\_\allowbreak from\_\allowbreak arbitrary\_\allowbreak stack}):
\begin{equation}
  \neg\begin{aligned}[t]
    \Bigl(&\mathrm{evalStack}(\mathrm{inv},\mathrm{cap},w_{00}) = \mathrm{evalQueue}(\mathsf{false},3,w_{00}) \\
          &\land\; \mathrm{evalStack}(\mathrm{inv},\mathrm{cap},w_{01}) = \mathrm{evalQueue}(\mathsf{false},3,w_{01})\Bigr).
  \end{aligned}
\end{equation}
\end{theorem}

Теорема~\ref{thm:h2_queue_vs_stack} закрывает сертификацию перехода на уровень очереди $L_4$ без постулированных таблиц истинности и без искусственного ограничения глубины стека $\mathrm{cap}$.

\subsection{Вычислимый детектор мнимых расширений сенсора}

Пусть базовый порт наблюдения задаёт отклик $\Pi_{\mathrm{base}}:Q\to\Omega\to\mathcal{O}_1$, а кандидат на расширенный порт возвращает пару $\Pi_{\mathrm{ext}}:Q\to\Omega\to\mathcal{O}_1\times\mathcal{O}_2$. Расширение является \emph{мнимым} ($\mathrm{IsSpuriousSemantic}$), если на конечном базисе проб $Q$ и активных миров $W$ первая компонента $\Pi_{\mathrm{ext}}$ совпадает с $\Pi_{\mathrm{base}}$, а любой сплит $\Pi_{\mathrm{ext}}(q,w_1)\ne\Pi_{\mathrm{ext}}(q,w_2)$ происходит только на тех парах миров $(w_1,w_2)\in W\times W$, которые \emph{уже} различаются базовым портом $\Pi_{\mathrm{base}}(q,w_1)\ne\Pi_{\mathrm{base}}(q,w_2)$.

Вместо постулирования универсального квантора в структуре данных мы определяем чистую булеву функцию $\mathrm{checkSpuriousFin}(W,Q,\Pi_{\mathrm{base}},\Pi_{\mathrm{ext}})\in\mathrm{Bool}$, перебирающую $Q\times W\times W$.

\begin{theorem}[Звучность вычислимого детектора мнимых сенсоров]
\label{thm:h3_spurious}
Для любых конечных списков миров $W$ и проб $Q$ с разрешимым равенством откликов (\texttt{H3\_checkSpuriousFin\_sound}):
\begin{equation}
  \mathrm{checkSpuriousFin}(W,Q,\Pi_{\mathrm{base}},\Pi_{\mathrm{ext}}) = \mathsf{true} \implies \mathrm{IsSpuriousSemantic}(W,Q,\Pi_{\mathrm{base}},\Pi_{\mathrm{ext}}).
\end{equation}
На трёхуровневом проверочном стенде детектор безошибочно отбраковывает детерминированный хеш-дубликат чётности ($\mathsf{true}$) и пропускает датчик скрытого старшего бита ($\mathsf{false}$, \texttt{H3\_detector\_\allowbreak separates\_\allowbreak spurious\_\allowbreak from\_\allowbreak genuine}).
\end{theorem}

\section{Единый контроллер с пятимерным бюджетом и разделением привилегий}

Объединим теперь все допустимые операции наблюдателя в единое индуктивное отношение перехода $\mathrm{LicensedStep}(\mathcal{T}, s, s')$, параметризованное замороженным положительным расписанием тарифов $\mathcal{T} = (c_q, c_s, c_p, c_v, c_m) > 0$ с условием амортизации памяти $c_m < c_s$. Состояние контроллера $s\in\mathrm{UnifiedState}$ хранит число активных миров $|W_s|$, число активных гипотез $|H_s|$, ранг текущего уровня языка $\mathrm{tier}(s)\in\{L_0,\dots,L_4\}$, ранг сенсорного порта $\mathrm{portRank}(s)$, интервенционный ранг $\mathrm{intRank}(s)$, счётчик накопленных сертификатов $\mathrm{knowledgeCount}(s)$, а также 5-мерный вектор израсходованных ресурсов $\mathrm{spent}(s)\in\mathbb{N}^5$ и векторный потолок бюджета $B\in\mathbb{N}^5$.

Контроллер допускает ровно шесть конструкторов перехода, каждый из которых защищён собственным структурным сертификатом и покомпонентной проверкой векторного бюджета $\mathrm{spent}(s') \le B$:
\begin{enumerate}
  \item \texttt{probe}: информационная проба, требующая строгого сжатия множества миров ($0 < |W_{s'}| < |W_s|$) со списанием тарифа $c_q$.
  \item \texttt{escalateLang}: повышение уровня языка $\mathrm{tier}(s) < \mathrm{tier}(s')$, требующее структурного сигнала неопровержимости $\mathrm{hasStructCert}(\mathrm{sig})=\mathsf{true}$ со списанием $c_s$.
  \item \texttt{escalatePort}: расширение сенсорного порта $\mathrm{portRank}(s')=\mathrm{portRank}(s)+1$, требующее как структурного сертификата препятствия, так и отрицательного вердикта детектора мнимости ($\mathrm{isSpurious}=\mathsf{false}$ из Теоремы~\ref{thm:h3_spurious}) со списанием $c_p$.
  \item \texttt{intervene}: активное хирургическое вмешательство, требующее сертификата $\mathrm{Gate15Cert}$ (пассивная коллизия наблюдения + ортогональность контрольной маски $K(do(I)(x))=K(x)$ + ненулевой коммутатор $[T_L,T_R]\ne 0$) со списанием $c_v$.
  \item \texttt{stochasticReject}: стохастическая отбраковка гипотез по последовательному $e$-значению, требующая сертификата $\mathrm{Gate16Cert}$ (явная положительность знаменателя $0<\mathrm{den}$ и полной массы $0<\mathrm{totalMass}$, наличие равномерной границы зазора $\mathrm{hasGapBound}=\mathsf{true}$ и пересечение порога $\mathrm{eValueNum}\ge \alpha^{-1}\cdot\mathrm{eValueDen}$) со списанием $c_q$.
  \item \texttt{recoverMemory}: восстановление ранее сертифицированного представления из ревизионной памяти при выполнении предусловия переноса $\mathrm{precondOk}=\mathsf{true}$ со списанием сниженного тарифа $c_m < c_s$.
\end{enumerate}

\begin{theorem}[Звучность, разделение привилегий и строгий прогресс единого контроллера]
\label{thm:h4_h7_controller}
Для любого положительного расписания тарифов $\mathcal{T}$ и любого лицензированного перехода $\mathrm{LicensedStep}(\mathcal{T}, s, s')$:
\begin{enumerate}
  \item \textbf{Инвариант векторного бюджета (\texttt{H4\_licensed\_\allowbreak step\_\allowbreak preserves\_\allowbreak vector\_\allowbreak budget}):} по всем пяти координатам $(c_q,c_s,c_p,c_v,c_m)$ выполняется $\mathrm{spent}(s') \le B$.
  \item \textbf{Разделение привилегий (\texttt{H4\_privilege\_\allowbreak separation\_\allowbreak blocks\_\allowbreak timeout\_\allowbreak escalation}):} сигнал ресурсного таймаута $\mathrm{resourceTimeout}$ без структурного свидетельства $\mathrm{certifiedUnsat}$ никогда не проходит страж $\mathrm{hasStructCert}=\mathsf{true}$ и не может вызвать \texttt{escalateLang} или \texttt{escalatePort}.
  \item \textbf{Блокировка мнимых портов и невалидных сертификатов (\texttt{H5\_spurious\_\allowbreak port\_\allowbreak blocked}, \texttt{H6\_gate15\_\allowbreak and\_\allowbreak gate16\_\allowbreak certificates\_\allowbreak enforced}):} при $\mathrm{isSpurious}=\mathsf{true}$ шлюз расширения порта возвращает $\mathsf{none}$, а шаги \texttt{intervene} и \texttt{stochasticReject} гарантируют выполнение всех алгебраических и вероятностных стражей.
  \item \textbf{Строгий монотонный прогресс (\texttt{H7\_licensed\_\allowbreak step\_\allowbreak strict\_\allowbreak epistemic\_\allowbreak progress}):} каждый шаг строго увеличивает суммарный индекс эпистемического прогресса $\mathcal{P}(s) = \mathrm{tier}(s) + \mathrm{portRank}(s) + \mathrm{intRank}(s) + \mathrm{knowledgeCount}(s)$, исключая вырожденные стратегии <<всегда НЕИЗВЕСТНО>> и циклы по мнимым сенсорам.
\end{enumerate}
\end{theorem}

\section{Алгоритмический синтез разделяющих признаков и типизированная память переноса}

\subsection{Алгоритмический оператор выбора разделяющей оси}

Пусть кандидаты $d\in M$ из разных предметных областей снабжены $k$-битовыми структурными сигнатурами признаков $\mathrm{features}(d):\mathrm{Fin}\,k\to\mathrm{Bool}$. Определим три оператора без участия человека в составлении меню проб: проекцию признака $\mathrm{projFeature}(a,d)=d.\mathrm{features}(a)$, проверку осевого расщепления
\begin{equation}
  \mathrm{isAxialSplit}(M, a) :\iff \bigl(\exists d_0\in M:\ \neg\mathrm{projFeature}(a,d_0)\bigr) \;\land\; \bigl(\exists d_1\in M:\ \mathrm{projFeature}(a,d_1)\bigr),
\end{equation}
и оператор заморозки пробы $\mathrm{freezeAxisQuery}(a,M,b) = \{\,d\in M \mid \mathrm{projFeature}(a,d)=b\,\}$ с гарантированным минимаксным выигрышем
\begin{equation}
  \mathrm{minimaxGain}(M,a) := |M| - \max_{b\in\{\mathsf{false},\mathsf{true}\}}\bigl|\mathrm{freezeAxisQuery}(a,M,b)\bigr|.
\end{equation}

\begin{theorem}[Гарантированный минимаксный выигрыш алгоритмического синтеза пробы]
\label{thm:h8_metasynset}
Для любого конечного списка междоменных объектов $M$ и любой структурной оси $a\in\mathrm{Fin}\,k$:
\begin{enumerate}
  \item (\texttt{H8\_partition\_\allowbreak length\_\allowbreak conservation}) Сумма мощностей ветвей равна исходному списку: $|\mathrm{freezeAxisQuery}(a,M,\mathsf{false})| + |\mathrm{freezeAxisQuery}(a,M,\mathsf{true})| = |M|$.
  \item (\texttt{H8\_algorithmic\_\allowbreak metasynset\_\allowbreak minimax\_\allowbreak gain\_\allowbreak ge\_1}) Если $\mathrm{isAxialSplit}(M,a)=\mathsf{true}$, то обе ветви строго сжимают список кандидатов и $\mathrm{minimaxGain}(M,a)\ge 1$.
\end{enumerate}
\end{theorem}

\subsection{Типизированная память переноса и дихотомия фона в Rule~110}

Токен ревизионной памяти $\tau\in\mathrm{RevisionToken}$ сохраняет домен происхождения, ширину порта $\mathrm{portWidth}$, допустимое число внешних входов шва $\mathrm{requiredSeamInputs}$, границу остаточного зазора $\Delta(T)$, цену открытия с нуля $c_s$ и цену переноса $c_m$. При переходе к новой среде $E_{\mathrm{env}}$ функция $\mathrm{verifyTransferPrecond}(\tau,E_{\mathrm{env}})$ проверяет совпадение ширины порта, автономность порта на новом фоне ($\mathrm{isPortAutonomous}=\mathsf{true}$), непревышение числа внешних входов шва ($\mathrm{actualSeamInputs}\le\mathrm{requiredSeamInputs}$), допуск по зазору и неравенство $c_m < c_s$.

Чтобы продемонстрировать, почему проверка физического фона среды критична для безопасности переноса, рассмотрим 4-клеточное окно наблюдения $W=\{3,4,5,6\}$ на 14-клеточном кольце клеточного автомата Rule~110 ($f_{110}(l,c,r)=c\oplus r\oplus cr\oplus lcr$):
\begin{itemize}
  \item В \textbf{модели с фиксированными соседями} значения $x_2=1$ и $x_7=0$ удерживаются на каждом такте по определению. Порт $W$ тогда обновляется замкнутым отображением $\mathrm{etherAutomaton}:\mathrm{Bit4}\to\mathrm{Bit4}$ без внешних входов. Это граничное условие модели, а не свойство свободного эфирного кольца.
  \item На \textbf{свободном 14-клеточном эфирном кольце} с начальной конфигурацией $\mathtt{00010011011111}$ (нумерация слева направо от 0) пара $(x_2,x_7)$ меняется. Прямой расчёт даёт следующую орбиту порта:
\[
\mathtt{1001}\to\mathtt{1011}\to\mathtt{1110}\to\mathtt{0010}\to
\mathtt{0110}\to\mathtt{1111}\to\mathtt{1000}\to\mathtt{1001}.
\]
  Семь различных портовых слов задают автономное табличное обновление только на этой орбите. В частности, свободное кольцо даёт переход $\mathtt{1011}\to\mathtt{1110}$, тогда как модель с фиксированными соседями даёт $\mathtt{1011}\to\mathtt{1111}$. Поэтому применение её детектора к этому штатному переходу свободного кольца даст ложный сигнал аномалии. Этот расчёт не является новой теоремой Lean и не устанавливает автономность при произвольных возмущениях.
  \item В \textbf{проверке на нулевом фоне} изменение правого соседа $x_7\in\{0,1\}$ при прочих нулевых входах меняет состояние клетки~6 уже за один такт. Это проверяемый пример неавтономности. Значение $\mathrm{actualSeamInputs}=4$ задано в записи среды теста переноса; эта проверка не доказывает минимальность числа входов шва.
\end{itemize}

\begin{theorem}[Звучность переноса памяти, потоковая амортизация и блокировка ложного переноса]
\label{thm:h9_h11_memory}
Для оператора восстановления токена $\mathrm{recoverToken}(\mathrm{mem}, E_{\mathrm{env}})$:
\begin{enumerate}
  \item (\texttt{H9\_recoverToken\_\allowbreak preserves\_\allowbreak provenance\_\allowbreak and\_\allowbreak gap}) Любой восстановленный токен $\tau$ принадлежит памяти $\mathrm{mem}$, сохраняет ширину порта и удовлетворяет всем ограничениям автономности, числа входов шва и остаточного зазора $\Delta(T)$.
  \item (\texttt{H10\_memory\_\allowbreak transfer\_\allowbreak strict\_\allowbreak stream\_\allowbreak amortization}) На любом потоке из $k\ge 2$ совместимых эпизодов суммарные затраты с памятью строго меньше затрат без памяти:
  \begin{equation}
    c_s + (k-1)c_m < k\cdot c_s.
  \end{equation}
  \item (\texttt{H11\_unsound\_\allowbreak transfer\_\allowbreak blocked\_\allowbreak on\_\allowbreak r110\_\allowbreak zero\_\allowbreak bg\_\allowbreak and\_\allowbreak d4\_\allowbreak shortcut}) Для заданной записи \texttt{envR110EtherRing} с $\mathrm{isPortAutonomous}=\mathsf{true}$ и $\mathrm{actualSeamInputs}=0$ проверка возвращает \texttt{some r110EtherPortToken}. Для двух других записей (нулевой фон Rule~110 и граф $D_4$ с ребром-шорткатом) она возвращает \texttt{none}. H11 проверяет результат на этих записях; он не выводит автономность свободного эфирного кольца.
\end{enumerate}
\end{theorem}

\section{Синтез эпистемической ревизии и горизонта ремонта}

Заключительный шаг связывает выбор эпистемического действия с физическим временем эволюционирующей среды $T:\Omega\to\Omega$. Пусть к дедлайну $t^*$ необходимо обеспечить целевой инвариант $\mathcal{I}:\Omega\to\mathrm{Bool}$, а локальный актуатор располагает набором вмешательств $A_{\mathrm{loc}}$. Как показано в динамической теории ремонта, множество физически достижимых до дедлайна вмешательств в момент $t\le t^*$ равно
\begin{equation}
  \RepairReach(W,t,t^*) = \bigl\{\, a\in A_{\mathrm{loc}} \;\big|\; \forall w\in W:\ \mathcal{I}\bigl(T^{t^*-t}(a(w))\bigr)=\mathsf{true} \,\bigr\}.
\end{equation}

В статическом контроллере при коллизии наблюдений на пассивном порту ($\Pi(w_0)=\Pi(w_1)$) естественным шагом является запуск пробы с задержкой $\Delta t>0$, разделяющей миры на эволюционировавших ветвях $T^{\Delta t}(W_{q,y})$. Однако в среде с конечной скоростью распространения сигнала ожидание $\Delta t$ тактов может привести к тому, что на всех ветвях $y$ множество $\RepairReach(T^{\Delta t}(W_{q,y}),\, t+\Delta t,\, t^*)$ окажется пустым.

Для непустого $W$ селектор \texttt{selectCausalEpistemicAction} проверяет, не истекает ли дедлайн за время ожидания и остаётся ли ремонт возможным на каждой достижимой ветви. Если ожидание не проходит эту проверку, а сейчас имеется общее вмешательство $a^*\in\RepairReach(W,t,t^*)$, селектор выбирает немедленный ремонт. Сам селектор не сравнивает информационный выигрыш или стоимость стратегий.

\begin{theorem}[Каузально-эпистемический перехват на 6-кольце Rule~110 и детекция аномалии субъекта]
\label{thm:h12_synthesis}
На 6-клеточном кольце Rule~110 для пары миров $W=\{\omega_0=(1,0,0,0,1,0),\ \omega_1=(1,0,0,1,1,0)\}$, дедлайна $t^*=3$, пассивного порта в клетке~$1$ с задержкой $\Delta t=2$, актуатора в клетке~$2$ и цели $w_0(3)=0$:
\begin{enumerate}
  \item (\texttt{H12\_causal\_\allowbreak epistemic\_\allowbreak horizon\_\allowbreak synthesis}) На такте $t=0$ миры неразличимы через пассивный порт ($\Pi_1(\omega_0)=\Pi_1(\omega_1)=\mathsf{false}$), а после ожидания $\Delta t=2$ полностью идентифицируются ($\Pi_1(T^2\omega_0)=\mathsf{false},\ \Pi_1(T^2\omega_1)=\mathsf{true}$); тем не менее на обеих ветвях на такте $t=2$ горизонт ремонта закрыт ($\RepairReach=\varnothing$), тогда как на $t=0$ действие \texttt{set2To1} спасает оба мира, и объединённый контроллер выбирает $\texttt{immediateSurgicalRepair\ set2To1}$.
  \item (\texttt{H12\_rule110\_\allowbreak algebraic\_\allowbreak xor\_\allowbreak gate}, \texttt{H12\_temporal\_\allowbreak subject\_\allowbreak period2\_\allowbreak and\_\allowbreak anomaly\_\allowbreak trigger}) Локальное обновление Rule~110 при единичном левом соседе реализует точный алгебраический гейт $f_{110}(1,z,k)=z\oplus k$ темпорального субъекта, а в 4-битной модели с фиксированными соседями состояние $\mathtt{0111}$ возвращается в себя за 2~такта, а на штатном первом переходе флаг аномалии не срабатывает ($\mathrm{subjectDetectsAnomaly}=\mathsf{false}$), тогда как подача $\mathtt{1111}$ вместо ожидаемого следующего состояния поднимает флаг аномалии ($\mathsf{true}$), но этот флаг сам по себе не удостоверяет допустимость или успешность ремонта.
\end{enumerate}
\end{theorem}

\begin{remark}[Точный объём теорем $H_{12}$ и связь с \texttt{v1.5}]
\label{rem:h12_scope}
Теорема \texttt{H12\_temporal\_\allowbreak subject\_\allowbreak period2\_\allowbreak and\_\allowbreak anomaly\_\allowbreak trigger} формально доказывает возврат состояния 4-битной модели с фиксированными соседями $\mathtt{0111}$ за 2~такта ($\mathrm{etherAutomaton}^2(\mathtt{0111})=\mathtt{0111}$, этот результат не относится к свободной эфирной орбите) и мгновенное срабатывание флага аномалии на $\mathtt{1111}$. Теорема \texttt{H12\_causal\_\allowbreak epistemic\_\allowbreak horizon\_\allowbreak synthesis} формализует случай закрытия горизонта при дедлайне $t^*=3$ ($\Delta t=2$); сравнение с дедлайном $t^*=4$ (где ожидание сохраняет $\RepairReach\neq\varnothing$ и проба становится необходимой), общая эквивалентность $V_{\mathrm{dyn}}$ деревьям планов и параметрическая хиральная граница $3m\to 2m$ на бесконечной решётке $\mathbb{Z}$ вынесены в сопутствующий модуль \texttt{causal\_dynamic\_repair\_v1\_5\_fixed.lean}.
\end{remark}

\section{Машинная верификация в Lean~4.36.0, воспроизводимость и границы}

Все определения и все 27~теорем разделов~2--6 формализованы в едином самодостаточном файле \texttt{certified\_\allowbreak strategy\_\allowbreak selection\_\allowbreak v2\_6\_1.lean} (1253~строки, 59\,667~байт, чистый LF без BOM, 0~внешних импортов, без Mathlib, пространство имён \texttt{CertifiedStrategySelectionV26}; SHA-256 \texttt{4c76155e475a7964\ldots{}f9cf9c77}). Модуль скомпилирован и проверен ядром \texttt{Lean 4.36.0} (\texttt{arm64-apple-darwin24.6.0}, commit \texttt{611c117b}): \texttt{exit 0}, 0~ошибок, 0~предупреждений, 0~\texttt{sorryAx}, 0~\texttt{Lean.ofReduceBool} (\texttt{native\_decide} не используется). Полное распределение аксиом по протоколу \texttt{V2\_6\_1\_VERIFICATION.log} приведено в Таблице~\ref{tab:axioms_v26}.

\begin{table}[!htbp]
\centering
\caption{Фактический профиль \texttt{\#print axioms} для всех 27~теорем в \texttt{certified\_strategy\_selection\_v2\_6\_1.lean} (Lean~4.36.0, \texttt{exit 0}, 0~\texttt{sorryAx}, 0~\texttt{native\_decide}).}
\label{tab:axioms_v26}
\footnotesize
\begin{tabularx}{\linewidth}{@{}l >{\raggedright\arraybackslash}X@{}}
\toprule
\textbf{Профиль \texttt{\#print axioms}} & \textbf{Верифицированные теоремы (27/27)} \\
\midrule
Без аксиом (10) & \texttt{H1\_identifiable\_\allowbreak implies\_\allowbreak expressible}, \texttt{H1\_identifiable\_\allowbreak implies\_\allowbreak distinguishable}, \texttt{H2\_queue\_\allowbreak recognizes\_\allowbreak fifo\_\allowbreak copy}, \texttt{H2\_stack\_\allowbreak recognizes\_\allowbreak lifo\_\allowbreak palindrome}, \texttt{H3\_detector\_\allowbreak separates\_\allowbreak spurious\_\allowbreak from\_\allowbreak genuine}, \texttt{H8\_cross\_\allowbreak domain\_\allowbreak instance\_\allowbreak verified}, \texttt{H11\_unsound\_\allowbreak transfer\_\allowbreak blocked\_\allowbreak on\_\allowbreak r110\_\allowbreak zero\_\allowbreak bg\_\allowbreak and\_\allowbreak d4\_\allowbreak shortcut}, \texttt{H12\_causal\_\allowbreak epistemic\_\allowbreak horizon\_\allowbreak synthesis}, \texttt{H12\_rule110\_\allowbreak algebraic\_\allowbreak xor\_\allowbreak gate}, \texttt{H12\_temporal\_\allowbreak subject\_\allowbreak period2\_\allowbreak and\_\allowbreak anomaly\_\allowbreak trigger} \\
\texttt{[propext]} (6) & \texttt{H1\_E2\_distinguishable\_\allowbreak not\_\allowbreak expressible}, \texttt{H3\_checkSpuriousFin\_sound}, \texttt{H4\_licensed\_\allowbreak step\_\allowbreak preserves\_\allowbreak vector\_\allowbreak budget}, \texttt{H4\_privilege\_\allowbreak separation\_\allowbreak blocks\_\allowbreak timeout\_\allowbreak escalation}, \texttt{H5\_spurious\_\allowbreak port\_\allowbreak blocked}, \texttt{H6\_gate15\_\allowbreak and\_\allowbreak gate16\_\allowbreak certificates\_\allowbreak enforced} \\
\texttt{[propext, Quot.sound]} (7) & \texttt{H1\_E1\_expressible\_\allowbreak not\_\allowbreak distinguishable}, \texttt{H1\_E3\_expr\_\allowbreak and\_\allowbreak dist\_\allowbreak not\_\allowbreak identifiable}, \texttt{H1\_three\_\allowbreak coordinate\_\allowbreak strict\_\allowbreak separation}, \texttt{H2\_stack\_\allowbreak inv\_\allowbreak false\_\allowbreak rejects\_\allowbreak w01}, \texttt{H2\_stack\_\allowbreak inv\_\allowbreak true\_\allowbreak rejects\_\allowbreak w00}, \texttt{H2\_gate13\_\allowbreak queue\_\allowbreak separates\_\allowbreak from\_\allowbreak arbitrary\_\allowbreak stack}, \texttt{H8\_partition\_\allowbreak length\_\allowbreak conservation} \\
\texttt{[propext, Classical.choice, Quot.sound]} (4) & \texttt{H7\_licensed\_\allowbreak step\_\allowbreak strict\_\allowbreak epistemic\_\allowbreak progress}, \texttt{H8\_algorithmic\_\allowbreak metasynset\_\allowbreak minimax\_\allowbreak gain\_\allowbreak ge\_1}, \texttt{H9\_recoverToken\_\allowbreak preserves\_\allowbreak provenance\_\allowbreak and\_\allowbreak gap}, \texttt{H10\_memory\_\allowbreak transfer\_\allowbreak strict\_\allowbreak stream\_\allowbreak amortization} \\
\bottomrule
\end{tabularx}
\end{table}

Сопровождающий эталонный пакет на чистом stdlib Python~3 (\texttt{certified\_\allowbreak strategy\_\allowbreak selection\_\allowbreak v2\_6.py}, SHA-256 \texttt{3dbdf0e5\ldots}, с явной защитой $W\neq\varnothing$ в \texttt{repair\_reach}; \texttt{test\_v2\_6.py}, SHA-256 \texttt{963d89ff\ldots}) исполняет все 12~проверочных сценариев и негативные тесты fail-closed аудитора \texttt{S0} (отклонение поддельных сертификатов с $\mathrm{den}=0$, $\mathrm{isSpurious}=\mathsf{true}$, неортогональной маской или сигналом $\mathrm{resourceTimeout}$), формируя детерминированный отчёт \texttt{v2\_6\_witness\_report.json} (SHA-256 \texttt{2ddf09ec\ldots}).

\paragraph{Пределы.}
Доказанные теоремы относятся к конечному списку активных миров $W$ и конечному базису проб $Q$ в каждый дискретный момент принятия решения (за исключением Теоремы~\ref{thm:h2_queue_vs_stack}, квантифицированной по всем $\mathrm{cap}\in\mathbb{N}$, и Теорем~\ref{thm:h8_metasynset}--\ref{thm:h9_h11_memory}, универсально квантифицированных по длине списка кандидатов $|M|$ и длине потока эпизодов $k\ge 2$). В частности: (1)~Теорема~\ref{thm:h2_queue_vs_stack} отделяет очередь $\mathrm{evalQueue}$ от конкретного 1-параметрического стекового вычислителя $\mathrm{evalStack}(\mathrm{inv},\mathrm{cap})$ произвольной ёмкости, а не от произвольного магазинного автомата (PDA) с произвольным конечным управлением; (2)~Теорема~\ref{thm:h4_h7_controller} гарантирует покоординатное сохранение 5-мерного бюджета $\mathrm{spent}(s')\le B$ на шагах $\mathrm{LicensedStep}$, но не решает задачу многокритериальной Парето-оптимизации по несравнимым бюджетам; (3)~как отмечено в Замечании~\ref{rem:h12_scope}, $H_{12}$ в данном модуле фиксирует дедлайн $t^*=3$ и 2-шаговый цикл модели с фиксированными соседями из состояния $\mathtt{0111}$, тогда как контраст с $t^*=4$ и хиральные границы на $\mathbb{Z}$ находятся в \texttt{causal\_dynamic\_repair\_v1\_5\_fixed.lean}. Стохастическая динамика самой среды $T$ и непрерывные пространства наблюдений в данную формализацию не входят и составляют предмет следующего шага программы.

\begin{thebibliography}{9}
\small\setlength{\itemsep}{2pt}

\bibitem{cook2004}
M.\,Cook. Universality in Elementary Cellular Automata. \emph{Complex Systems}, 15(1):1--40, 2004.

\bibitem{martinez2006}
G.\,J.\,Mart\'inez, H.\,V.\,McIntosh, J.\,C.\,Seck-Tuoh-Mora. Gliders in Rule~110. \emph{International Journal of Unconventional Computing}, 2(1):1--49, 2006.

\bibitem{reiter1987}
R.\,Reiter. \href{https://doi.org/10.1016/0004-3702(87)90062-2}{A Theory of Diagnosis from First Principles}. \emph{Artificial Intelligence}, 32(1):57--95, 1987.

\bibitem{hyafil1976}
L.\,Hyafil, R.\,L.\,Rivest. \href{https://www.sciencedirect.com/science/article/abs/pii/0020019076900958}{Constructing Optimal Binary Decision Trees Is NP-Complete}. \emph{Information Processing Letters}, 5(1):15--17, 1976.

\bibitem{wolpert2008}
D.\,H.\,Wolpert. \href{https://doi.org/10.1016/j.physd.2008.03.040}{Physical Limits of Inference}. \emph{Physica D: Nonlinear Phenomena}, 237(9):1257--1281, 2008.

\bibitem{pearl2009}
J.\,Pearl. \emph{Causality: Models, Reasoning, and Inference}, 2nd ed. Cambridge University Press, 2009.

\end{thebibliography}

\end{document}
